%! Change in Linear and Exponential Functions — Unit 4, Lesson 4.2 — Homework
\documentclass[10pt]{article}
\usepackage{precalculus-article}
\usepackage{precalculus-boxes}
\usepackage{grading}
\gradingassignment{precalc-u04-l02-hw}
\gradingcourse{precalculus}
\newcommand{\TallMath}[1]{$\displaystyle #1\rule[-1.4em]{0pt}{3.2em}$}

\begin{document}

\pageheader{Unit 4, Lesson 4.2}{Homework}

\vspace{0.10in}

\begin{enumerate}[itemsep=16pt]

\item For each table, apply the equal-length interval test and identify the function type.
      Then write the formula.

      \begin{enumerate}[label=(\alph*), itemsep=12pt]
        \item
        \renewcommand{\arraystretch}{1.2}
        \begin{tabular}{c|c}
          $x$ & $y$ \\ \hline
          0 & 80 \\ 1 & 40 \\ 2 & 20 \\ 3 & 10
        \end{tabular}

        \vspace{0.06in}
        Differences: \smash{\gradedblank[id=hw1a1, type=short, check=text, points=1, answer={-40, -20, -10}, display={$-40, -20, -10$}]{4cm}{13pt}} \quad Ratios: \smash{\gradedblank[id=hw1a2, type=short, check=text, points=1, answer={0.5, 0.5, 0.5}, display={$0.5, 0.5, 0.5$}]{4cm}{13pt}}

        Type: \smash{\gradedblank[id=hw1a3, type=short, check=text, points=1, answer={exponential}]{3cm}{13pt}} \quad Formula: $f(x) = $ \smash{\gradedblank[id=hw1a4, type=short, points=2, answer={80 \cdot (0.5)^x}]{4cm}{13pt}}

        \item
        \begin{tabular}{c|c}
          $x$ & $y$ \\ \hline
          1 & 5 \\ 3 & 11 \\ 5 & 17 \\ 7 & 23
        \end{tabular}

        \vspace{0.06in}
        (Interval length = \smash{\gradedblank[id=hw1b1, type=short, points=1, answer={2}]{1.5cm}{13pt}}) \quad Differences: \smash{\gradedblank[id=hw1b2, type=short, check=text, points=1, answer={6, 6, 6}, display={$6, 6, 6$}]{4cm}{13pt}}

        Type: \smash{\gradedblank[id=hw1b3, type=short, check=text, points=1, answer={linear}]{3cm}{13pt}} \quad Formula: $f(x) = $ \smash{\gradedblank[id=hw1b4, type=short, points=2, answer={3x + 2}]{4cm}{13pt}}
      \end{enumerate}

\item Two points are given: $(3,\ 4)$ and $(7,\ 20)$. Construct a \textbf{linear} function
      through these points.

      \vspace{0.06in}
      Slope: $m = \dfrac{20 - 4}{7 - 3} = $ \smash{\gradedblank[id=hw2a, type=short, points=1, answer={4}]{2cm}{13pt}}

      \vspace{0.06in}
      Function: $f(x) = $ \smash{\gradedblank[id=hw2b, type=short, points=2, answer={4x - 8}]{6cm}{13pt}}

\item Two points are given: $(1,\ 8)$ and $(4,\ 1000)$. Construct an \textbf{exponential}
      function through these points.

      \vspace{0.06in}
      $b = \left(\dfrac{1000}{8}\right)^{1/(4-1)} = \left(\text{\smash{\gradedblank[id=hw3a, type=short, points=1, answer={125}]{2cm}{10pt}}}\right)^{1/3} = $ \smash{\gradedblank[id=hw3b, type=short, points=1, answer={5}]{2cm}{13pt}}

      \vspace{0.06in}
      Function: $f(x) = $ \smash{\gradedblank[id=hw3c, type=short, points=2, answer={1.6 \cdot 5^x}]{6cm}{13pt}}

\item A savings account starts with \$500.

      \begin{enumerate}[label=(\alph*), itemsep=8pt]
        \item \textbf{Plan A} adds \$50 every month. Write a linear function $A(t)$ for the balance
              after $t$ months.

              $A(t) = $ \smash{\gradedblank[id=hw4a, type=short, points=1, answer={500 + 50t}]{5cm}{13pt}}

        \item \textbf{Plan B} grows by 5\% each month. Write an exponential function $B(t)$ for
              the balance after $t$ months.

              $B(t) = $ \smash{\gradedblank[id=hw4b, type=short, points=1, answer={500 \cdot (1.05)^t}]{5cm}{13pt}}

        \item Calculate $A(12)$ and $B(12)$. Which plan gives a higher balance after 12 months?

              \vspace{0.06in}
              $A(12) = \$$ \smash{\gradedblank[id=hw4c1, type=short, points=1, answer={1100}]{3cm}{13pt}} \quad $B(12) \approx \$$ \smash{\gradedblank[id=hw4c2, type=short, check=numeric, tolerance=0.005, points=2, answer={897.93}]{3cm}{13pt}}

              \vspace{0.04in}
              Higher balance: \smash{\gradedblank[id=hw4c3, type=short, check=text, points=1, answer={Plan A}]{4cm}{13pt}}
      \end{enumerate}

\item \textbf{(AP-Style Multiple Choice)} Which of the following tables could represent an
      exponential function?

      \medskip
      \begin{enumerate}[label=(\Alph*), itemsep=4pt]
        \item
        \begin{tabular}{c|c}
          $x$ & $y$ \\ \hline 0 & 5 \\ 1 & 8 \\ 2 & 11 \\ 3 & 14
        \end{tabular}

        \item
        \begin{tabular}{c|c}
          $x$ & $y$ \\ \hline 0 & 3 \\ 1 & 6 \\ 2 & 12 \\ 3 & 24
        \end{tabular}

        \item
        \begin{tabular}{c|c}
          $x$ & $y$ \\ \hline 0 & 1 \\ 1 & 4 \\ 2 & 9 \\ 3 & 16
        \end{tabular}

        \item
        \begin{tabular}{c|c}
          $x$ & $y$ \\ \hline 0 & 0 \\ 1 & 2 \\ 2 & 4 \\ 3 & 8
        \end{tabular}
      \end{enumerate}

      \vspace{0.06in}
      \leavevmode\rlap{\textit{Letter:}}\phantom{Answer:}\spacefactor2000{} \smash{\gradedblank[id=hw5a, type=choice, points=1, answer=B]{1.5cm}{13pt}} \quad Explain: \writegap\\
      \writebox{\raisebox{-0.5pt}[0pt][0pt]{\gradedblank[id=hw5b, type=rubric, points=2,
  display={Table B: consecutive ratios are all 2 ($6/3 = 12/6 = 24/12$), so outputs change proportionally: exponential.}]{\linewidth}{20pt}}}\\

\item \textbf{(Extension)} Let $f(x) = ab^x$ pass through two points $(x_1, y_1)$ and
      $(x_2, y_2)$ with $y_1 \ne 0$.

      \begin{enumerate}[label=(\alph*), itemsep=8pt]
        \item Write the two equations $y_1 = ab^{x_1}$ and $y_2 = ab^{x_2}$.

        \item Divide the second equation by the first and solve for $b$.

              \vspace{0.06in}
              Show algebraically that $b = \left(\dfrac{y_2}{y_1}\right)^{1/(x_2 - x_1)}$.

              \vspace{0.8in}
      \end{enumerate}

\end{enumerate}

\end{document}
